\subsection{限制形式幂级数}

定义 1。若交换拓扑环 A 存在一个由 A 的理想组成的 0 的邻域基本系 B，则称 A 具有线性拓扑（并称其拓扑是线性的）。

注意，在这样的环中，\(\mathfrak{S} \in \mathcal{B}\) 中的理想是开且闭的（一般拓扑学，第三章，第 2 节，第 1 号，命题 4 的推论）。对于所有 \(\mathfrak{S} \in \mathcal{B}\)，商拓扑环 A/\(\mathfrak{S}\) 是离散的；对于 \(\mathfrak{S} \in \mathcal{B}\)、\(\mathfrak{S}' \in \mathcal{B}\)、\(\mathfrak{S}' \subset \mathfrak{S}\)，令

\[
h _ {\mathfrak {S S} ^ {\prime}} \colon \mathrm{A} / \mathfrak {S} ^ {\prime} \rightarrow \mathrm{A} / \mathfrak {S}
\]

是典范映射。我们知道（一般拓扑学，第三章，第 7 节，第 3 号），\((\mathbf{A} / \mathfrak{S}, h_{\mathfrak{S}\mathfrak{S}'})\) 是离散环的逆系统（指标集为 \(\mathcal{B}\)，按 \(\supset\) 排序且有向），其逆极限是完备的 Hausdorff 线性拓扑环 \(\tilde{\mathbf{A}}\)；此外（同上，命题 2），定义了严格态射 \(i: \mathbf{A} \to \tilde{\mathbf{A}}\)，其核是 \(\{0\}\) 在 \(\mathbf{A}\) 中的闭包，像在 \(\tilde{\mathbf{A}}\) 中处处稠密，因此 \(\tilde{\mathbf{A}}\) 典范地识别为 \(\mathbf{A}\) 的 Hausdorff 完备化。

定义 2。给定交换拓扑环 A，形式幂级数

\[
\mathbf {T} = \sum_ {(n _ {i})} c _ {n _ {1} n _ {2} \dots n _ {p}} \mathbf {X} _ {1} ^ {n _ {1}} \mathbf {X} _ {2} ^ {n _ {2}} \dots \mathbf {X} _ {p} ^ {n _ {p}}
\]

若属于环 \(\mathbf{A}[[\mathbf{X}_1,\ldots ,\mathbf{X}_p]]\)，并且对于 A 中 0 的每个邻域 \(\mathbf{V}\)，系数 \(c_{n_1n_2\dots n_p}\) 中不属于 \(\mathbf{V}\) 的只有有限多个，则称其为限制的（换言之，族 \((c_{n_1n_2\dots n_p})\) 关于 \(\mathbf{N}^p\) 的有限子集补集滤子在 A 中趋于 0）。

若 A 具有线性拓扑，则 \(A[[X_{1},\ldots,X_{p}]]\) 中的限制形式幂级数构成 \(A[[X_{1},\ldots,X_{p}]]\) 的子环，记为 \(A\{X_{1},\ldots,X_{p}\}\)：因为若 \(T=\sum_{(n_{i})}c_{n_{1}\ldots n_{p}}X_{1}^{n_{1}}\ldots X_{p}^{n_{p}}\)、\(T'=\sum_{(n_{i})}c_{n_{1}\ldots n_{p}}^{'}X_{1}^{n_{1}}\ldots X_{p}^{n_{p}}\) 是两个限制形式幂级数，且 \(\mathfrak{S}\) 是 A 中 0 的一个邻域并且是 A 的理想，则存在整数 \(m\)，使得 \(c_{n_1\ldots n_p} \in \mathfrak{S}\) 且 \(c_{n_1\ldots n_p}' \in \mathfrak{S}\) 对每个指标组 \((n_1, \ldots, n_p)\) 都成立，只要至少有一个 \(n_k \geqslant m\)，其中 \(k\) 为指标；现在，若

\[
\mathrm{T} ^ {\prime \prime} = \mathrm{TT} ^ {\prime} = \sum_ {(n _ {i})} c _ {n _ {1} \dots n _ {p}} ^ {\prime \prime} \mathrm{X} _ {1} ^ {n _ {1}} \dots \mathrm{X} _ {p} ^ {n _ {p}},
\]

于是 \(c_{n_1 \ldots n_p}'' = \sum c_{r_1 \ldots r_p} c_{s_1 \ldots s_p}'\)，其中指标组 \((r_k), (s_k)\) 满足 \(r_k + s_k = n_k\)（\(1 \leqslant k \leqslant p\)）。我们得到，若 \(n_k \geqslant 2m\)，则 \(r_k \geqslant m\) 或 \(s_k \geqslant m\)，因而由于 \(\mathfrak{S}\) 是理想，有 \(c_{n_1 \ldots n_p}'' \in \mathfrak{S}\)，只要 \(n_k \geqslant 2m\) 对至少一个 \(k\) 成立，这就证明了我们的断言。此外，限制形式幂级数的每个导数 \(\partial T / \partial X_i\)（\(1 \leqslant i \leqslant p\)）都是限制的，这立即由定义以及邻域 \(\mathfrak{S} \in \mathcal{B}\) 是 A 的加法子群这一事实得出。

若 A 是离散的，则限制形式幂级数环就是多项式环 \(A[X_{1},\ldots,X_{p}]\)。

以下始终假设 A 具有线性拓扑，并令 \(\mathcal{B}\) 为 A 中由 A 的理想组成的 0 的邻域基本系；对于所有 \(\mathfrak{S} \in \mathcal{B}\)，令 \(p_{\mathfrak{S}}: A \to A/\mathfrak{S}\) 为典范同态。根据定义，对于每个限制形式幂级数 \(T \in A\{X_{1}, \ldots, X_{p}\}\)，

\[
\bar {p} _ {\Im} (\mathrm{T}) \in (\mathrm{A} / \Im) [ \mathrm{X} _ {1}, \dots , \mathrm{X} _ {p} ].
\]

显然

\[
((\mathbf {A} / \mathfrak {S}) [ \mathbf {X _ {1}}, \dots , \mathbf {X _ {p}} ], \bar {h} _ {\mathfrak {S} \mathfrak {S} ^ {\prime}})
\]

是环的逆系统（相对于有向指标集 \(\mathcal{B}\)），而 \((\bar{p}_{3})\) 是同态的逆系统 \(\mathrm{A}\{\mathrm{X}_1,\dots ,\mathrm{X}_p\} \to (\mathrm{A} / \Im)[\mathrm{X}_1,\dots ,\mathrm{X}_p]\)；由于每个多项式都是限制形式幂级数，\(\bar{p}_{3}\) 是满射；其核 \(\mathrm{N}_{3}\) 是 \(\mathrm{A}\{\mathrm{X}_1,\dots ,\mathrm{X}_p\}\) 的理想，由所有系数都属于 \(\Im\) 的限制形式幂级数组成；赋予 \(\mathrm{A}\{\mathrm{X}_1,\dots ,\mathrm{X}_p\}\) 以（线性）拓扑，使 \(\mathrm{N}_{3}\)（其中 \(\Im \in \mathcal{B}\)）构成 0 的邻域基本系（这个拓扑显然只取决于 A 上的拓扑）。于是由一般拓扑学第三章第 7 节第 3 号命题 2 可知

\[
\pi = \varprojlim_ {\mathfrak {S}} \bar {p} _ {\mathfrak {S}}: \mathrm{A} \{\mathrm{X} _ {1}, \dots , \mathrm{X} _ {p} \} \rightarrow \varprojlim_ {\mathfrak {S}} (\mathrm{A} / \mathfrak {S}) [ \mathrm{X} _ {1}, \dots , \mathrm{X} _ {p} ]\tag{3}
\]

是严格态射，其核是 \(\{0\}\) 在 \(\mathbf{A}\{\mathbf{X}_1,\ldots ,\mathbf{X}_p\}\) 中的闭包，且其像在

\[
\mathrm{A} ^ {\prime} = \lim _ {\leftarrow \mathfrak {S}} (\mathrm{A} / \mathfrak {S}) [ \mathrm{X} _ {1}, \dots , \mathrm{X} _ {p} ].
\]

命题 3。若具有线性拓扑的交换环 A 是 Hausdorff 且完备的，则典范同态 \(\pi\) 是拓扑环同构。

对所有 \((n_1, \ldots, n_p) \in \mathbf{N}^p\) 以及所有 \(\mathfrak{S} \in \mathcal{B}\)，令 \(\phi_{n_1 \ldots n_p}^{\mathfrak{S}}\) 为映射 \((\mathrm{A} / \mathfrak{S})[\mathrm{X}_1, \ldots, \mathrm{X}_p] \to \mathrm{A} / \mathfrak{S}\)，它将每个多项式映到该多项式中 \(\mathrm{X}_1^{n_1} \ldots \mathrm{X}_p^{n_p}\) 的系数；显然，\(\phi_{n_1 \ldots n_p}^{\mathfrak{S}}\) 构成 \((\mathrm{A} / \mathfrak{S})\)-模同态的逆系统（相对于有序集 \(\mathcal{B}\)），并且根据假设，\(\mathbf{A}\) 典范地识别为 \(\varprojlim_{\mathfrak{S}} (\mathrm{A} / \mathfrak{S})\)，所以 \(\phi_{n_1 \ldots n_p} = \varprojlim_{\mathfrak{S}} \phi_{n_1 \ldots n_p}^{\mathfrak{S}}\) 是从 \(\mathrm{A}'\) 到 \(\mathbf{A}\) 的连续 A-同态。对于每个元素 \(S = (S_{\mathfrak{S}})_{\mathfrak{S} \in \mathcal{B}}\)，属于 \(\mathrm{A}'\)，我们将看到，形式幂级数 \(T = \sum_{(n_i)} \phi_{n_1 \ldots n_p}(S) X_1^{n_1} \ldots X_p^{n_p}\) 是限制的，并满足 \(\pi(T) = S\)。对于所有 \(\mathfrak{S} \in \mathcal{B}\) 以及所有 \(\mathfrak{S}' \in \mathcal{B}\) 满足 \(\mathfrak{S}' \subset \mathfrak{S}\)，若关系 \(\phi_{n_1 \ldots n_p}^{\mathfrak{S}}(S_{\mathfrak{S}}) = 0\)，则

\[
\phi_ {n _ {1} \dots n _ {p}} ^ {\mathfrak {S} ^ {\prime}} (\mathrm{S} _ {\mathfrak {S} ^ {\prime}}) \in \mathfrak {S} / \mathfrak {S} ^ {\prime};
\]

由于 \(S_{\mathfrak{S}}\) 是多项式，我们看到，\(\phi_{n_{1}\ldots n_{p}}(S) \in \mathfrak{S}\)，除了有限多个 \((n_{1}, \ldots, n_{p})\) 之外；这些有限多个指标组满足 \(\phi_{n_{1}\ldots n_{p}}^{\mathfrak{S}}(S_{\mathfrak{S}}) \neq 0\)。这证明了第一个断言；第二个断言由定义得出。由于 A 是 Hausdorff 的，\(N_{\mathfrak{S}}\) 的交化为 0，因而 \(\pi\) 是双射；由于 \(\pi\) 是严格态射，这就完成了证明。

命题 4。设 A、B 为两个具有线性拓扑的交换环，B 是 Hausdorff 且完备的，\(u: \mathbf{A} \to \mathbf{B}\) 是连续同态。对于 B 中的每个元素族 \(\mathbf{b} = (b_i)_{1 \leq i \leq p}\)，存在唯一的连续同态

\[
\tilde {u} \colon \mathrm{A} \{\mathrm{X} _ {1}, \dots , \mathrm{X} _ {p} \} \rightarrow \mathrm{B}
\]

使得 \(\tilde{u}(a) = u(a)\) 对所有 \(a \in \mathbf{A}\) 都成立，并且 \(\tilde{u}(\mathbf{X}_i) = b_i\) 对 \(1 \leqslant i \leqslant p\) 都成立。

存在唯一同态 \(v \colon A[X_1, \ldots, X_p] \to B\)，使得 \(v(a) = u(a)\) 对 \(a \in A\) 都成立，并且 \(v(X_i) = b_i\) 对 \(1 \leqslant i \leqslant p\) 都成立。此外，若 \(\mathfrak{H}\) 是 \(B\) 中的一个 0 的邻域且为理想，则 \(\bar{u}^{-1}(\mathfrak{H}) = \mathfrak{I}\) 是 \(A\) 的一个理想且为 0 的邻域，并且对于每个多项式 \(P \in N_{\mathfrak{S}}\)，显然 \(v(P) \in \mathfrak{H}\)，所以 \(v\) 连续。由于 \(A[X_1, \ldots, X_p]\) 在 \(A\{X_1, \ldots, X_p\}\) 中稠密，\(\tilde{u}\) 的存在性和唯一性由一般拓扑学第三章第 3 节第 3 号命题 5 以及恒等式延拓原理得出。

在 \(\mathbf{A} = \mathbf{B}\) 且 \(u\) 为恒等映射的特殊情形，我们把 \(f(b_{1},\ldots ,b_{p})\) 或 \(f(\mathbf{b})\) 作为 \(\tilde{u} (f)\) 的值，其中 \(f\in \mathrm{A}\{\mathbf{X}_1,\dots ,\mathbf{X}_p\}\)。

\textbf{注}\par

\begin{enumerate}
\def\labelenumi{(\arabic{enumi})}
\item
  命题 4 证明了：对于一个假定为 Hausdorff 且完备的环 A 中的每个闭理想 \(\mathfrak{a}\)，关系 \(b_{i} \in \mathfrak{a}\) 对 \(1 \leqslant i \leqslant p\) 成立时，都有 \(f(b_{1}, \ldots, b_{p}) \in \mathfrak{a}\)，其中每个限制形式幂级数 \(f \in \mathrm{A}\{\mathrm{X}_{1}, \ldots, \mathrm{X}_{p}\}\) 均可取。
\item
  设 \(\mathbf{A}\) 具有线性拓扑；令 \(r\) 为满足
\end{enumerate}

\(1 \leqslant r \leqslant p\) 的整数，并赋予环 \(\mathrm{A}\{\mathrm{X}_1, \ldots, \mathrm{X}_r\}\) 以上述拓扑。则拓扑环 \(\mathrm{A}\{\mathrm{X}_1, \ldots, \mathrm{X}_p\}\) 可识别为限制形式幂级数环

\[
(\mathrm{A} \{\mathrm{X} _ {1}, \dots , \mathrm{X} _ {r} \}) \{\mathrm{X} _ {r + 1}, \dots , \mathrm{X} _ {p} \}
\]

这立即由定义得出。

\begin{enumerate}
\def\labelenumi{(\arabic{enumi})}
\setcounter{enumi}{2}
\tightlist
\item
  采用注 2 的记号，再假设 A 是 Hausdorff 且完备的，并将每个限制形式幂级数 \(f \in \mathrm{A}\{\mathbf{X}_1, \ldots, \mathbf{X}_p\}\) 写成如下形式：
\end{enumerate}

\[
f = \sum_ {(n _ {i})} c _ {n _ {r + 1} \dots n _ {p}} (\mathrm{X} _ {1}, \dots , \mathrm{X} _ {r}) \mathrm{X} _ {r + 1} ^ {n _ {r + 1}} \dots \mathrm{X} _ {p} ^ {n _ {p}}
\]

其中 \(c_{n_{r+1}\ldots n_p}\) 是限制形式幂级数。对于 A 中元素组成的每个组 \(\mathbf{x} = (x_1, \ldots, x_r)\)，令

\[
b _ {n _ {r + 1} \dots n _ {p}} = c _ {n _ {r + 1} \dots n _ {p}} (x _ {1}, \dots , x _ {p}).
\]

由注 1 立即可见，\(\sum_{(n_i)} b_{n_r+1 \ldots n_p} X_{r+1}^{n_r+1} \ldots X_p^{n_p}\) 是限制形式幂级数，记为 \(f(x_1, \ldots, x_r, X_{r+1}, \ldots, X_p)\)；称其为用 \(x_i\) 代替 \(X_i\)（对 \(1 \leqslant i \leqslant r\)）于形式幂级数 \(f\) 中所得到的级数。

